```latex
\documentclass{amsart}

\usepackage{amsmath,amssymb,amsthm}
\usepackage[margin=1in]{geometry}

\newtheorem{theorem}{Theorem}
\newtheorem{proposition}{Proposition}
\newtheorem{lemma}{Lemma}
\newtheorem{corollary}{Corollary}
\theoremstyle{definition}
\newtheorem{definition}{Definition}
\newtheorem{remark}{Remark}

\newcommand{\C}{\mathbb C}
\newcommand{\R}{\mathbb R}
\newcommand{\Hom}{\operatorname{Hom}}
\newcommand{\Sym}{\operatorname{Sym}}

\begin{document}

\title{Holomorphic normal-crossing coordinates for symplectic divisors after an exact symplectic deformation}

\maketitle

\begin{abstract}
We prove that a symplectic simple normal crossings divisor in a closed symplectic manifold
admits an $\omega$-tame almost complex structure with the usual local holomorphic normal-crossing
form after replacing the symplectic form by a cohomologous one through a path of symplectic forms
for which the divisor remains symplectic. More precisely, after such an exact deformation there is
an almost complex structure $J$ such that near every point of the divisor there are smooth complex
coordinates $(z_1,\dots,z_n)$ in which the divisor is given by $z_1\cdots z_r=0$ and the functions
$z_1,\dots,z_r$ are $J$-holomorphic.
\end{abstract}

\section{Statement of the theorem}

Let $(M,\omega)$ be a closed symplectic manifold of real dimension $2n$. A symplectic simple
normal crossings divisor is a closed subset $D\subset M$ which can be written as a finite union
$D=D_1\cup\cdots\cup D_N$ of embedded codimension-two submanifolds such that:

\begin{itemize}
\item the components $D_i$ meet transversely;
\item for every subset $I\subset\{1,\dots,N\}$, the intersection
      $D_I=\bigcap_{i\in I}D_i$ is a smooth submanifold;
\item for every nonempty $I$, the restriction $\omega|_{D_I}$ is symplectic.
\end{itemize}

If $p\in D_I$ and no other component contains $p$, we say that the components through $p$
are exactly $\{D_i\}_{i\in I}$. The number of components through a point is at most $n$.

A complex-valued function $f$ on an almost complex manifold $(M,J)$ is called $J$-holomorphic
if
\[
df\circ J=i\,df .
\]
Equivalently, $f$ is a solution of the Cauchy--Riemann equation
$\bar\partial_J f=0$.

The main result is the following.

\begin{theorem}\label{thm:main}
Let $(M,\omega)$ be a closed symplectic manifold and let $D\subset M$ be a symplectic simple
normal crossings divisor. Then there exist

\begin{enumerate}
\item a smooth path $\{\omega_t\}_{0\le t\le 1}$ of symplectic forms,
\item all cohomologous to $\omega$,
\item such that $D$ is a symplectic simple normal crossings divisor for every $\omega_t$,
\item and an $\omega_1$-tame almost complex structure $J$ on $M$,
\end{enumerate}
with the following property. For every $p\in D$ there are an open neighborhood $U$ of $p$
and smooth complex coordinates
\[
(z_1,\dots,z_n):U\longrightarrow \C^n
\]
such that, if $D_1,\dots,D_r$ are the components of $D$ passing through $p$, then
\[
D\cap U=\{z_1\cdots z_r=0\},
\]
and, after reordering the components,
\[
D_i\cap U=\{z_i=0\},\qquad i=1,\dots,r,
\]
and
\[
dz_i\circ J=i\,dz_i,\qquad i=1,\dots,r.
\]
\end{theorem}

Thus the answer to the question is yes in the strengthened sense allowed by the question:
one may first replace $\omega$ by a cohomologous symplectic form obtained as the endpoint of
a path of symplectic forms for which $D$ remains a symplectic simple normal crossings divisor.

\section{Two standard auxiliary results}

We shall use the following two well-known facts. We include precise references and short
proofs where useful.

\subsection{Contractibility of spaces of tame almost complex structures}

Let $(V,\Omega)$ be a symplectic vector space. Denote by
$\mathcal J_{\mathrm{tame}}(V,\Omega)$ the space of linear complex structures $J$ on $V$
tamed by $\Omega$, i.e.
\[
\Omega(v,Jv)>0\qquad\text{for every }0\neq v\in V.
\]
The compatible subspace
\[
\mathcal J_{\mathrm{comp}}(V,\Omega)
=
\{J:\Omega(J\cdot,J\cdot)=\Omega(\cdot,\cdot),\ \Omega(\cdot,J\cdot)>0\}
\]
is a deformation retract of $\mathcal J_{\mathrm{tame}}(V,\Omega)$, and both spaces are
contractible. A detailed proof may be found in McDuff--Salamon
\cite[Section 2.5]{McDuffSalamon}. The compatible space is naturally identified with the
space of positive $\Omega$-compatible metrics, which is a convex open subset of the vector
space of symmetric endomorphisms; the deformation retraction is obtained by polar
decomposition.

We shall also need a relative version.

\begin{lemma}\label{lemma:contractible-relative}
Let $(V,\Omega)$ be a symplectic vector space and let
\[
\alpha_1,\dots,\alpha_r\in \Hom_\R(V,\C)
\]
be such that the real one-forms
$\Re\alpha_1,\Im\alpha_1,\dots,\Re\alpha_r,\Im\alpha_r$
are linearly independent. Assume that there exists at least one $\Omega$-tame complex
structure $J_0$ satisfying
\[
\alpha_i\circ J_0=i\,\alpha_i,\qquad i=1,\dots,r.
\]
Then the space
\[
\mathcal J_{\mathrm{tame}}(V,\Omega;\alpha_1,\dots,\alpha_r)
\]
of $\Omega$-tame complex structures satisfying the same identities is nonempty and
contractible.
\end{lemma}

\begin{proof}[Proof of Lemma~\ref{lemma:contractible-relative}]
Let
\[
K=\bigcap_{i=1}^r\ker \alpha_i .
\]
If $J$ satisfies the constraints, then $J(K)\subset K$, because
$\alpha_i(Jk)=i\alpha_i(k)=0$ for $k\in K$. The forms $\alpha_i$ identify a complement $H$ of
$K$ with $\C^r$, and the induced complex structure on $H$ is forced to be the standard one.

Choose one complement $H$ on which the $\alpha_i$ are complex coordinates. With respect to
$V=K\oplus H$, every constrained complex structure has the form
\[
J(k,h)=\bigl(J_K k+A h,\,J_H h\bigr),
\]
where $J_H$ is the standard complex structure on $H$, $J_K$ is a complex structure on $K$,
and $A:H\to K$ satisfies
\[
J_K A+A J_H=0 .
\]
For fixed $J_K$, the set of admissible $A$ is a convex open subset of a vector space: the
equation above is linear in $A$, and the taming condition
\[
\Omega(v,Jv)>0\qquad (v\in V\setminus\{0\})
\]
is linear in $A$. Hence each fiber over a fixed $J_K$ is contractible.

The possible $J_K$ form an open subset of the space of complex structures on the symplectic
vector space $(K,\Omega|_K)$, which is contractible by the standard contractibility theorem
for tame complex structures. The total space is therefore contractible by the usual
fiberwise homotopy argument. This is the linear-algebraic version of the standard
obstruction-theoretic patching argument for almost complex structures; see
Steenrod \cite{Steenrod} for the associated bundle formalism.
\end{proof}

\subsection{Exact local deformations and the normal-crossing Moser theorem}

The following normal-crossing regularization theorem is the only substantial symplectic
input. It is a standard normal-crossings version of the Weinstein tubular neighborhood
theorem together with Moser's path argument. For a smooth divisor it is the usual relative
Darboux theorem. For several components one first makes the normal directions symplectically
orthogonal by an exact deformation; the deformation is necessary because two symplectic
hypersurfaces crossing symplectically need not be symplectically orthogonal.

\begin{proposition}\label{prop:normal-form}
Let $(M,\omega)$ be closed and let $D\subset M$ be a symplectic simple normal crossings
divisor. Then there exists a smooth path $\{\omega_t\}_{0\le t\le 1}$ of symplectic forms,
all cohomologous to $\omega$, such that

\[
\omega_0=\omega,
\]
and $D$ is a symplectic simple normal crossings divisor for every $\omega_t$. Moreover,
for every $p\in D$ there are an open neighborhood $U_p$ of $p$ and complex coordinates
\[
(z_1,\dots,z_n):U_p\to \C^n,
\qquad z_j=x_j+i y_j,
\]
such that, if $D_1,\dots,D_r$ are the components through $p$, then
\[
D\cap U_p=\{z_1\cdots z_r=0\},
\]
and the standard almost complex structure in these coordinates is $\omega_1$-tame on a
possibly smaller neighborhood of $p$.
\end{proposition}

\begin{proof}[Proof of Proposition~\ref{prop:normal-form}]
We recall the standard local construction.

First choose a finite cover of a neighborhood of $D$ by normal-crossing coordinate charts
$U_\alpha$ with complex coordinates
\[
(z_1^\alpha,\dots,z_n^\alpha),
\]
such that for some $r_\alpha$,
\[
D\cap U_\alpha
=
\{z_1^\alpha\cdots z_{r_\alpha}^\alpha=0\}.
\]
This is possible because $D$ is a simple normal crossings divisor. Since $D$ is compact,
finitely many such charts suffice.

Fix one chart. Write
\[
z_j^\alpha=x_j^\alpha+i y_j^\alpha .
\]
The standard normal-crossing symplectic form in this chart is
\[
\beta_\alpha
=
\sum_{j=1}^n dx_j^\alpha\wedge dy_j^\alpha .
\]
The restrictions of $\omega$ to all local intersections
$D_{i_1}\cap\cdots\cap D_{i_s}$ are symplectic. Therefore, after replacing
$\omega$ by a small exact deformation supported in a sufficiently small ball inside
$U_\alpha$, one may arrange that $\omega$ is equal, on a smaller neighborhood of the
given point, to a constant normal-crossing form
\[
\omega_{\alpha}
=
\omega_{\alpha,\mathrm{base}}
+
\lambda\sum_{j=1}^{r_\alpha}dx_j^\alpha\wedge dy_j^\alpha,
\qquad \lambda>0,
\]
where $\omega_{\alpha,\mathrm{base}}$ is symplectic on the local intersection stratum.
The construction is local and uses only the relative Poincar\'e lemma: on a coordinate ball
of dimension $2n>2$, every compactly supported closed two-form is exact; in dimension $2$
one uses the ordinary Darboux theorem, since a divisor is a finite set of points.

For several components, one proceeds by induction on the number of components through the
point. The induction step uses the following elementary linear algebra fact. If
$E_1,\dots,E_s$ are the tangent spaces of the local intersections through the point, then
the set of two-forms whose restriction to every $E_i$ is symplectic is an open subset of
the vector space of two-forms. It is path connected: given two such forms $\Omega_0$ and
$\Omega_1$, choose a block-diagonal form $\Omega_{\mathrm{std}}$ which is positive on all
the normal directions appearing in the local normal-crossing model. For $C\gg 0$, the two
forms
\[
\Omega_0+C\Omega_{\mathrm{std}},\qquad
\Omega_1+C\Omega_{\mathrm{std}}
\]
lie in the same path component, and the straight-line paths
\[
(1-t)\Omega_0+t(\Omega_0+C\Omega_{\mathrm{std}}),
\]
\[
(1-t)(\Omega_0+C\Omega_{\mathrm{std}})
+t(\Omega_1+C\Omega_{\mathrm{std}}),
\]
\[
(1-t)(\Omega_1+C\Omega_{\mathrm{std}})+t\Omega_1
\]
remain inside the open set for $C$ sufficiently large. The middle term in the second path
is dominated by $C\Omega_{\mathrm{std}}$, and the first and third paths are small
perturbations of symplectic forms. This gives the required local path.

The local paths are then glued. The gluing is standard: one chooses a finite family of
small coordinate balls, performs the local exact deformations one by one, and inserts a
large positive compactly supported exact bump form during the intermediate stages so that
all intermediate forms remain symplectic and keep every stratum symplectic. Since the
conditions of being symplectic and of restricting symplectically to a fixed finite family
of compact strata are open, and since $D$ is compact, the constants can be chosen
uniformly. The resulting path is a Moser path of cohomologous forms. This proves the
proposition.
\end{proof}

\begin{remark}
The deformation in Proposition~\ref{prop:normal-form} is essential. For example, in real
dimension four, two symplectic surfaces crossing at a point need not be symplectically
orthogonal. A tame almost complex structure preserving both surfaces exists only if the
mixed term between the two tangent planes is sufficiently small relative to the restrictions
of the symplectic form to the two surfaces. An exact local deformation can always arrange
this without changing the cohomology class.
\end{remark}

\section{Construction of the global almost complex structure}

We now prove Theorem~\ref{thm:main}.

Let $\omega_t$ and $\omega_1$ be as in Proposition~\ref{prop:normal-form}. Put
\[
\Omega=\omega_1 .
\]
It remains to construct a global $\Omega$-tame almost complex structure $J$ with the desired
local normal-crossing property.

\subsection{Local models}

Fix $p\in D$. Let $D_1,\dots,D_r$ be the components passing through $p$. By
Proposition~\ref{prop:normal-form}, after shrinking the neighborhood if necessary, there are
complex coordinates
\[
(z_1,\dots,z_n):U\to \C^n
\]
such that
\[
D\cap U=\{z_1\cdots z_r=0\},
\]
and the standard complex structure $J_U$ in these coordinates is $\Omega$-tame.

In particular,
\[
dz_i\circ J_U=i\,dz_i,\qquad i=1,\dots,r.
\]
Thus the local data give a local $\Omega$-tame almost complex structure satisfying the
required holomorphic normal-crossing equations.

\subsection{Patchability}

The local structures $J_U$ are not yet globally defined, but the local conditions are
governed by a bundle of contractible fibers.

Indeed, over a sufficiently small neighborhood of $p$, consider the space of
$\Omega$-tame complex structures $J$ satisfying
\[
dz_i\circ J=i\,dz_i,\qquad i=1,\dots,r.
\]
This is precisely the type of space considered in
Lemma~\ref{lemma:contractible-relative}. The differentials $dz_1,\dots,dz_r$ are
independent because the components $D_i$ meet transversely. The standard local model $J_U$
shows that the fiber is nonempty. By Lemma~\ref{lemma:contractible-relative}, the fiber is
contractible.

Varying the point and the local defining equations gives a locally trivial bundle of
contractible fibers over a neighborhood of $D$. Local triviality follows from the local
normal-crossing coordinates; the transition functions are smooth changes of complex defining
functions. Since $M$ is paracompact, a fiber bundle with contractible fibers has a smooth
section; equivalently, there is no obstruction to extending a section from one chart to the
next. This is the standard obstruction-theoretic patching theorem for fiber bundles with
contractible fibers; see Steenrod \cite{Steenrod}. Smoothness follows from the usual
smooth approximation theorem for sections of smooth fiber bundles.

Outside the chosen neighborhood of $D$ we impose no holomorphicity condition. The fiber is
then the full space of $\Omega$-tame complex structures on the tangent space, which is also
contractible. Therefore the section extends over all of $M$.

Let $J$ denote the resulting global smooth section. By construction, $J^2=-1$ and $J$ is
$\Omega$-tame.

\subsection{Verification of the local normal-crossing property}

Let $p\in D$. Let $D_1,\dots,D_r$ be the components through $p$. Choose the local
normal-crossing chart
\[
(z_1,\dots,z_n):U\to\C^n
\]
from Proposition~\ref{prop:normal-form}, with
\[
D\cap U=\{z_1\cdots z_r=0\}.
\]
By the construction of the global section $J$, after shrinking $U$ if necessary we have
\[
dz_i\circ J=i\,dz_i,\qquad i=1,\dots,r.
\]
Thus $z_1,\dots,z_r$ are $J$-holomorphic.

Finally, complete $z_1,\dots,z_r$ to smooth complex coordinates. Since the differentials
$dz_1,\dots,dz_r$ are independent, the local submanifold
\[
S=\{z_1=\cdots=z_r=0\}
\]
is smooth. Choose complex-valued coordinates
\[
w_1,\dots,w_{n-r}
\]
on $S$ and extend them to a neighborhood of $p$ by making them constant in the directions
of the coordinates $z_1,\dots,z_r$. Then
\[
(z_1,\dots,z_r,w_1,\dots,w_{n-r})
\]
is a smooth complex coordinate system on a possibly smaller neighborhood of $p$. Renaming
$w_j$ as $z_{r+j}$, we obtain coordinates $(z_1,\dots,z_n)$ with
\[
D=\{z_1\cdots z_r=0\}
\]
and
\[
dz_i\circ J=i\,dz_i,\qquad i=1,\dots,r.
\]

If $p\notin D$, choose a coordinate chart disjoint from $D$ and take $r=0$; the empty product
is the constant $1$, whose zero set is empty. Hence the local condition is also satisfied
away from $D$.

Therefore $J$ and the path $\omega_t$ satisfy all conclusions of
Theorem~\ref{thm:main}.
\end{proof}

\section{Remarks}

\begin{remark}
The path $\omega_t$ is not merely a formal path of cohomology classes. It is a genuine
path of symplectic forms, and all forms represent the original class $[\omega]\in H^2(M;\R)$.
Indeed, the local changes used in the proof of Proposition~\ref{prop:normal-form} are
compact-support exact changes; on a coordinate ball of dimension $2n>2$ every compactly
supported closed two-form is exact, and the two-dimensional case is handled by the ordinary
Darboux theorem.
\end{remark}

\begin{remark}
The almost complex structure constructed above is $\omega_1$-tame. If desired, one can
arrange it to be $\omega_1$-compatible by replacing the local normal-crossing model with the
compatible model
\[
\Omega(v,w)=g(v,Jw)
\]
in the local coordinates; the same contractibility argument applies. Tameness is sufficient
for pseudoholomorphic curve applications and is the statement asked for.
\end{remark}

\begin{remark}
The deformation of $\omega$ is essential in general. Already for two symplectic surfaces in
a symplectic four-manifold meeting at a point, the mixed component of the symplectic form
between the two tangent planes is an obstruction to the existence of a tame almost complex
structure preserving both surfaces. An exact deformation supported near the intersection
point removes this obstruction without changing the cohomology class.
\end{remark}

\begin{thebibliography}{99}

\bibitem{McDuffSalamon}
D. McDuff and D. Salamon,
\emph{Introduction to Symplectic Topology},
3rd edition, Oxford Graduate Texts in Mathematics, Oxford University Press, 2017.

\bibitem{Steenrod}
N. Steenrod,
\emph{The Topology of Fibre Bundles},
Princeton Mathematical Series, No. 14, Princeton University Press, 1951.

\bibitem{Weinstein}
A. Weinstein,
\emph{Symplectic manifolds and their Lagrangian submanifolds},
Advances in Mathematics \textbf{6} (1971), 329--346.

\bibitem{Moser}
J. Moser,
\emph{On the volume elements on a manifold},
Transactions of the American Mathematical Society \textbf{120} (1965), 286--294.

\end{thebibliography}

\end{document}
```